\chapter{اللقاء الأربعون: آية الدين والكتابة والإشهاد وحكمة التشريع}

\section{مقدمة}
السلام عليكم ورحمة الله وبركاته. الحمد لله رب العالمين، وأشهد أن لا إله إلا الله وحده لا شريك له، وأشهد أن محمداً عبده ورسوله، صلى الله عليه وعلى آله وصحبه وسلم.

دخل الشيخ في هذا اللقاء إلى قوله تعالى: \begin{ayah}يَا أَيُّهَا الَّذِينَ آمَنُوا إِذَا تَدَايَنْتُمْ بِدَيْنٍ إِلَى أَجَلٍ مُسَمًّى فَاكْتُبُوهُ\end{ayah}، وهي آية الدين، أطول آية في كتاب الله، وأكثرها تفصيلاً في حفظ الحقوق وصيانة المعاملات.

وفي هذا المجلس يتجلى فقه \rawi{الطبري} في أمرين جليلين: جمع صور التداين كلها تحت خطاب الآية، ثم بيان أن أوامر الشريعة في هذا الباب ليست صوراً شكلية، بل هي أسباب رحمة وصيانة واحتياط للناس في أموالهم وحقوقهم.

\separator

\section{عموم التداين وأن الآية لا تختص بالقرض}
\isnad{قال أبو جعفر:} معنى قوله تعالى: \begin{ayah}إِذَا تَدَايَنْتُمْ بِدَيْنٍ\end{ayah} إذا تعاطيتم معاملة ثبت بها مال في الذمة إلى أجل معلوم، سواء كان ذلك بيعاً بثمن مؤجل، أو سلماً، أو قرضاً، أو غير ذلك من صور المعاملات التي يبقى فيها الحق إلى وقت مؤجل.

ونبّه الشيخ هنا إلى أن بعض الناس يقصر الآية على باب القرض، مع أن \rawi{الطبري} يوسع مدلولها ليشمل كل مال ثبت لأحد على أحد بسبب معاملة مؤجلة. وذكر أن \rawi{ابن عباس} رضي الله عنهما خص سبب النزول بالسلم، لكن ذلك لا يمنع عموم الحكم في سائر الصور الداخلة في معنى التداين.

ثم وقف الشيخ عند دقة التعبير في قوله: \begin{ayah}بِدَيْنٍ\end{ayah}، وأنه ليس حشواً ولا تأكيداً مجرداً، بل جاء لرفع احتمال معنى آخر في لفظ التداين عند العرب، فإنهم قد يستعملونه في المجازاة أيضاً، فجاءت الآية لتعيين المراد هنا: التداين المالي الذي يحتاج إلى ضبط وكتابة.

\begin{fawaid}[ليست كل زيادة في اللفظ تأكيداً]
من أخطاء التفسير أن يقال في كل لفظ زائد ظاهراً: إنه للتأكيد. وكثير من هذه الألفاظ إنما جاء للبيان والفصل ورفع الاحتمال، كما في قوله تعالى: \begin{ayah}إِذَا تَدَايَنْتُمْ بِدَيْنٍ\end{ayah}.
\end{fawaid}

\section{فاكتبوه: بين الوجوب والندب}
عند قوله تعالى: \begin{ayah}فَاكْتُبُوهُ\end{ayah} عرض \rawi{الطبري} خلاف أهل العلم: هل الأمر هنا للوجوب أو للندب؟ ثم مال إلى أن الأصل في الأمر الوجوب حتى تقوم حجة صارفة، وأن الله أوجب على المتداينين صيانة هذا الحق بالكتابة ما دام ذلك ميسوراً.

وذكر الشيخ ما روي عن بعض السلف في التشديد في هذا الباب، حتى روي عن بعضهم أن الرجل قد يكون مظلوماً في ضياع حقه، ثم لا يستجاب له إذا كان قد فرط في الأسباب التي أمر الله بها، فلم يكتب ولم يشهد مع القدرة على ذلك.

ثم استخرج الشيخ من ذلك أصلاً كبيراً: أن كثيراً من الناس يظن أن الاحتياط الذي أمرت به الشريعة في الأموال والحقوق معناه سوء الظن، وليس كذلك، بل هو من كمال الرعاية للنفوس وللأموال. فقد ينسى الدائن، وقد ينسى المدين، وقد يموت أحدهما، وقد يقع النزاع بعد صفاء، فجعل الله الكتابة حصناً للحق قبل وقوع الخصومة.

\begin{fawaid}[الاحتياط للحقوق من محاسن الشريعة لا من سوء الظن]
كتابة الدين ليست اتهاماً للناس بقدر ما هي صيانة لهم من النسيان والموت والتغير والنزاع، ولهذا كانت من تمام العدل والرحمة بين المؤمنين.
\end{fawaid}

\section{تضييق باب النسخ في آية الدين}
ذكر \rawi{الطبري} قول من قال إن الأمر بالكتابة نسخ بقوله تعالى: \begin{ayah}فَإِنْ أَمِنَ بَعْضُكُمْ بَعْضًا\end{ayah}، ثم رده رداً نفيساً، وبيّن أن هذه الآية ليست ناسخة، لأن النسخ لا يكون إلا إذا امتنع الجمع بين الحكمين، أما هنا فالجمع ممكن.

فحمل \rawi{الطبري} آية الأمانة على حال تعذر الكتابة أو الكاتب، أو على حال مخصوصة لا تنافي بقاء أصل الأمر بالكتابة عند الإمكان. ومن هنا أعاد الشيخ تقرير أصل متكرر في هذه الدروس: أن \rawi{الطبري} يضيق باب النسخ غاية التضييق، ولا يقبل إبطال حكم ثابت إلا بحجة قاطعة لا تقبل التردد.

\begin{fawaid}[كلما أمكن الجمع سقطت دعوى النسخ]
ليس كل انتقال من حال إلى حال في القرآن نسخاً، بل متى أمكن حمل كل آية على موضعها وظرفها مع بقاء العمل بالجميع كان ذلك أولى من إهدار بعض النصوص باسم النسخ.
\end{fawaid}

\section{وليكتب بينكم كاتب بالعدل}
في قوله تعالى: \begin{ayah}وَلْيَكْتُبْ بَيْنَكُمْ كَاتِبٌ بِالْعَدْلِ\end{ayah}، بيّن \rawi{الطبري} أن المراد كاتب يكتب بالحق والإنصاف، لا يميل مع أحد الطرفين، ولا يترك حقاً، ولا يزيد باطلاً، ولا يصوغ العبارة على وجه يفتح باب النزاع أو التحيل.

ونبّه الشيخ إلى أن العدل هنا لا يقتصر على حسن النية، بل يدخل فيه أيضاً حسن المعرفة بصيغ المعاملات، ودقة العبارة، ومراعاة العرف الجاري بين الناس، حتى تكون الوثيقة قاطعة للنزاع لا مولدة له.

ولهذا كانت الشريعة في هذا الباب سابقة إلى ما يحتاجه الناس في كل زمان: من ضبط الأجل، وتعيين المال، وإزالة الاحتمال، وسد أبواب التنازع قبل وقوعه.

\section{ولا يأب كاتب أن يكتب كما علمه الله}
فسر \rawi{الطبري} هذه الآية بأن الكاتب لا ينبغي له أن يمتنع من الكتابة إذا احتيج إليه، ما دام قادراً غير متضرر، لأن الله هو الذي علمه هذه النعمة وخصه بها دون كثير من الناس.

ووقف الشيخ مع هذه اللفتة وقفة تربوية عامة، فذكر أن كل من خصه الله بنعمة أو مهارة أو علم أو صنعة يحتاجها الناس فشكر هذه النعمة أن يبذل منها، وأن يدخل بها في معنى قوله تعالى: \begin{ayah}وَمِمَّا رَزَقْنَاهُمْ يُنْفِقُونَ\end{ayah}. فالعلم رزق، والبيان رزق، والكتابة رزق، كما أن المال رزق.

\begin{fawaid}[كل من خصه الله بشيء فشكره أن يبذل منه]
ليست الصدقة في المال وحده، بل من شكر نعمة الله على العبد أن ينفع الناس بما علمه الله وأقدره عليه، ومن ذلك الكتابة والبيان وسائر المنافع التي يحتاج إليها الخلق.
\end{fawaid}

\separator

\section{أقسط عند الله وأقوم للشهادة وأدنى ألا ترتابوا}
عند قوله تعالى: \begin{ayah}ذَلِكُمْ أَقْسَطُ عِنْدَ اللَّهِ وَأَقْوَمُ لِلشَّهَادَةِ وَأَدْنَى أَلَّا تَرْتَابُوا\end{ayah}، بسط الشيخ المعنى الذي ذكره \rawi{الطبري}: أن الكتابة أعدل في ميزان الله، وأثبت للشهادة، وأقرب إلى رفع الشك والاختلاف.

وهنا أكثر الشيخ من بيان حكمة التشريع، وأن من أجمل ما يقوم به أهل العلم أن لا يقتصروا على نقل الأحكام المجردة، بل يبينوا للناس ما فيها من العدل والمصلحة. فكم من حقوق ضاعت في واقع الناس لأنهم فرطوا في الكتابة والإشهاد، ثم لم يعرفوا أن ما ضاع إنما ضاع بتركهم ما أمر الله به.

ومن أبلغ ما قاله في هذا الموضع أن المؤمن كلما ازداد أخذاً بأحكام الشريعة ازداد انتفاعاً بها في دنياه قبل آخرته، وهذا باب واسع في الطلاق، والبيوع، والجوار، والحدود، وسائر الأبواب.

\begin{fawaid}[بيان حكمة التشريع من أعظم أبواب التعليم]
إذا رأى الناس آثار أحكام الله في حفظ الدماء والأموال والبيوت والحقوق ازدادوا يقيناً بعدله وحكمته، وسهل عليهم الانقياد له والانشراح به.
\end{fawaid}

\section{إلا أن تكون تجارة حاضرة}
جاء الاستثناء في قوله تعالى: \begin{ayah}إِلَّا أَنْ تَكُونَ تِجَارَةً حَاضِرَةً تُدِيرُونَهَا بَيْنَكُمْ\end{ayah}، فبيّن \rawi{الطبري} أن المراد به التجارة الناجزة التي يتقابض فيها الطرفان الحقين في المجلس، فلا يبقى دين مؤجل يحتاج إلى توثيق.

ولهذا رخص الله في ترك الكتابة هنا، لأن موجبها قد زال بتقابض العوضين. ومع ذلك لم يسقط أصل الإرشاد إلى الإشهاد في البيع، لأن ضياع الحق قد يقع فيه أيضاً من جهة أخرى.

وتناول الشيخ في هذا الموضع خلاف القراءة والإعراب في \begin{ayah}تِجَارَةً حَاضِرَةً\end{ayah}، وبيّن أن \rawi{الطبري} رجح ما ظهر له من جهة العربية والسياق، لا على وجه رد القراءة الثابتة إذا صحت، بل على حسب ما بلغه من ثبوتها. واستثمر ذلك في التنبيه على الفرق بين ترجيح معنى في القراءات، وبين رد القراءة المتواترة نفسها.

\section{وأشهدوا إذا تبايعتم}
رجح \rawi{الطبري} أن الإشهاد في البيع أيضاً على الأصل من الوجوب أو اللزوم ما لم تقم حجة صريحة على صرفه، وبيّن أن الحكمة في ذلك لا تختص بباب الدين المؤجل.

فالبائع قد يجحد البيع، والمشتري قد يجحد الثمن أو أصل المعاملة، وقد يضيع حق أحدهما إذا خلا البيع من توثيق أو شهود. ومن هنا كانت الشريعة تحتاط للطرفين معاً، ولا تجعل التوثيق حماية لطرف دون طرف.

\begin{fawaid}[الشريعة تحتاط للطرفين لا لطرف واحد]
الكتابة والإشهاد في أبواب المعاملات ليست لمصلحة الدائن وحده ولا البائع وحده، بل هي شبكة أمان للطرفين، لأن الخصومة قد تأتي من أي جهة.
\end{fawaid}

\section{ولا يضار كاتب ولا شهيد}
عرض \rawi{الطبري} في تفسير هذه الآية أقوالاً، ثم رجح أن النهي متوجه إلى أصحاب المعاملة: فلا يضار الكاتب ولا الشهيد بأن يلجآ إلى ما فيه مشقة أو تعسف أو تعطيل لهما عن مصالحهما بغير حق.

وبيّن الشيخ أن هذا الترجيح يقوى من جهة السياق؛ لأن الخطاب في أكثر الآية متجه إلى أصحاب الحقوق والمتعاملين، فحمل هذا الموضع على النسق نفسه أولى من قطعه إلى معنى آخر بعيد.

ثم استخرج من ذلك قاعدة في الترجيح: أن توجيه اللفظ إلى ما يوافق سياق الآية ونظمها أولى من نقله إلى معنى منقطع عنها، ما دام اللسان يحتمله.

\begin{fawaid}[السياق من أقوى المرجحات في تعيين المخاطَب]
إذا احتمل الضمير أو الأمر أو النهي أكثر من جهة، كان حمله على من جرى عليهم الخطاب في الآيات قبلَه وبعدَه أولى من صرفه إلى غيرهم بلا قرينة قوية.
\end{fawaid}

\section{واتقوا الله ويعلمكم الله}
عند قوله تعالى: \begin{ayah}وَاتَّقُوا اللَّهَ وَيُعَلِّمُكُمُ اللَّهُ\end{ayah}، بيّن الشيخ أن المعنى ليس على ما يتوهمه بعض الناس من تعليق التعليم على مجرد دعوى التقوى على وجه غامض، بل المراد أن الله يبين لعباده الأحكام ويعلمهم ما ينفعهم، وعليهم أن يتقوه بالعمل بهذا البيان.

فالسياق هنا سياق أحكام ومعاملات وتفصيلات، لا سياق إشارات مجردة. ومن اتقى الله في هذا الباب بأن أخذ بالأسباب المشروعة، وانتفع بتعليم الله، كان ذلك من تمام الفقه والعمل.

\separator

خلاصة هذا اللقاء أن آية الدين مثال باهر على كمال الشريعة في الجمع بين العدل والرحمة والاحتياط؛ فهي لا تترك الناس لعاداتهم وتقلباتهم في الأموال، بل تضع لهم نظاماً دقيقاً في الكتابة والإشهاد والعدل وصيانة الوثائق والشهود. وفيها أيضاً درس منهجي راسخ في تضييق دعوى النسخ، وفي أن كثيراً من أوامر الله إنما شرعت رحمة بالخلق قبل أن تكون مجرد تكاليف عليهم.
